\section{Scope, Motivation, and Pipeline Placement}
\label{sec:c1_scope}

\subsection{Why Component 1 exists}
Hard-threshold pruning of retrieved evidence (for example, ``keep only if entailment score is above a cutoff'')
causes evidence starvation in multi-hop fact verification.
Component 1 replaces hard pruning with a reversible, auditable control layer:
\begin{itemize}
\item score every retrieved evidence item with an NLI-based premise validator,
\item partition evidence into Active/Suspended/Counter sets instead of deleting items,
\item compute sufficiency and safety metrics before downstream reasoning,
\item log all critical artifacts in machine-readable JSONL.
\end{itemize}

\subsection{What Component 1 is responsible for}
\textbf{In scope (implemented in code):}
\begin{itemize}
\item Evidence representation and deterministic verbalization from triples/sentences.
\item Premise Validation (PV) scoring with transformer-based NLI.
\item SQLite cache keyed by claim/evidence/model/verbalizer settings.
\item Evidence State Manager (ESM) partitioning into $\Active$, $\Susp$, and $\Counter$.
\item Sufficiency diagnostics: $\ESIprod$, $\ESIgeom$, $\CR@k$, $\RPI@k$, neutral dominance.
\item Structured logging for claim-level and pair-level outputs.
\end{itemize}

\textbf{Out of scope (explicitly not implemented in Component 1):}
\begin{itemize}
\item Graph reasoner/verifier architecture and training.
\item Multi-step recovery controller or backtracking loop.
\item Final claim classification logic beyond PV diagnostics and A/S/C assignment.
\end{itemize}

\subsection{Pipeline position and handoff semantics}
\begin{figure}[H]
\centering
\begin{adjustbox}{max width=\textwidth}
\begin{tikzpicture}[
  node distance=7mm and 8mm,
  block/.style={draw,rounded corners,align=center,text width=3.45cm,minimum height=1.3cm,fill=gray!8},
  arrow/.style={-Latex,thick},
  dashedarrow/.style={-Latex,thick,dashed}
]
\node[block] (claim) {FactKG claim row\\\texttt{Sentence, Label, Entity\_set}};
\node[block, right=of claim] (retrieved) {Retrieved subgraph row\\\texttt{walked["walkable"]}};
\node[block, right=of retrieved] (pv) {PV scoring\\$\Ent, \Con, \Neu, \Rel, \Pol$};
\node[block, right=of pv] (esm) {ESM partition\\$\Pool \rightarrow \Active/\Susp/\Counter$};
\node[block, below=of pv] (metrics) {Metrics + logs\\\texttt{claims.jsonl, pairs.jsonl}};
\node[block, right=of metrics] (c2) {Component 2 input\\soft masks + diagnostics};

\draw[arrow] (claim) -- (retrieved);
\draw[arrow] (retrieved) -- (pv);
\draw[arrow] (pv) -- (esm);
\draw[arrow] (esm.south) -- (metrics.north);
\draw[arrow] (metrics) -- (c2);

\draw[dashedarrow] (esm.east) ++(0.2,0.3) -- ++(1.2,0) node[midway,above,font=\scriptsize]{\Susp retained};
\draw[dashedarrow] (esm.east) ++(0.2,-0.3) -- ++(1.2,0) node[midway,below,font=\scriptsize]{\Counter retained};
\end{tikzpicture}
\end{adjustbox}
\caption{Page-fitted Component 1 flow. Active evidence is consumed immediately; Suspended and Counter are retained for diagnostics and future recovery logic.}
\label{fig:c1_pipeline}
\end{figure}

\subsection{Core notation used across the document}
\begin{table}[H]
\centering
\begin{tabular}{@{}p{0.20\linewidth}p{0.74\linewidth}@{}}
\toprule
Symbol & Meaning \\
\midrule
$\Claim$ & Claim text (hypothesis in NLI). \\
$\Pool$ & Full retrieved evidence pool for a claim. \\
$\Active,\Susp,\Counter$ & Active, Suspended, and Counter evidence sets (disjoint partition of $\Pool$). \\
$\Ent, \Con, \Neu$ & Entailment, contradiction, and neutral probabilities from PV. \\
$\Rel(e)$ & Relevance mass: $\Ent(e)+\Con(e)=1-\Neu(e)$. \\
$\Pol(e)$ & Polarity: $\Ent(e)-\Con(e)$. \\
$\ESIprod, \ESIgeom$ & Product and geometric-mean Evidence Sufficiency Index variants. \\
$\CR@k$ & Counter Retention at rank $k$. \\
$\RPI@k$ & Recovery Potential Index at rank $k$. \\
\bottomrule
\end{tabular}
\caption{Compact notation reference for Component 1.}
\label{tab:c1_notation}
\end{table}
